\documentclass[10pt,a4paper]{amsart}
\usepackage[margin=19mm]{geometry}
\usepackage{amsmath,amssymb,mathtools}
\usepackage[scr=rsfs,cal=euler]{mathalfa}
\usepackage{xfrac}
\usepackage[unicode,hidelinks,bookmarks=false]{hyperref}
\newcommand{\ie}{i.e.\ }
\newcommand{\cf}{cf.\ }
\newcommand{\resp}{respectively\ }
\newcommand{\Subsection}{\subsection}
\providecommand{\nobreakdash}{\nobreak\hbox{-}}
\setlength{\emergencystretch}{2em}
\multlinegap0pt
\usepackage{amssymb}
\newtheorem{lemma}{Lemma}
\newtheorem{theorem}{Theorem}
\title{Hamiltonian reduction from particular integrals}
\author{R. Azuaje, A. M. Escobar-Ruiz, I. Gutierrez-Sagredo}
\date{Source: arXiv:2607.07057v2 (CC BY 4.0)}
\begin{document}
\maketitle

\noindent\emph{Source and licence.} Hamiltonian reduction from particular integrals. Version arXiv:2607.07057v2 (CC BY 4.0), distributed by arXiv under the Creative Commons Attribution 4.0 International licence, http://creativecommons.org/licenses/by/4.0/, as recorded in the arXiv licence metadata for this exact version and shown on the abstract page https://arxiv.org/abs/2607.07057v2. Full licence notice: \texttt{licenses/arxiv-2607.07057-CC-BY-4.0.txt}.\par\smallskip

\noindent\textit{Editorial scope.} Counted unit: the article's theorem maintheorem (source0.tex lines 315--363). The article's complete original proof of it is delivered with it (source0.tex lines 364--413, one \texttt{proof} environment of 50 lines). No mathematics is written, repaired or re-derived: the statement and the proof are the article's own lines, in the article's own notation, and no retained line is substituted or re-flowed.  Retained uncounted as the source-local context the proof needs: lines 194--207.  Nothing was omitted from any retained span, so the package carries no omission record.  No retained line contains a citation, so the wrapper carries no bibliography and no bibliography entry.  No retained line contains a figure environment, an image-inclusion command or drawing code, so no image is delivered and the package declares no asset dependency.  Editorial formatting only, in addition: an AMS \texttt{amsart} preamble that restates the article's own declarations and macros used by the retained lines, namely the environments \texttt{lemma}, \texttt{theorem} on separate counters, together with the standard packages those lines need.  The retained environments print in this document as Lemma 1, Theorem 1 in order of appearance, whereas the article prints the same items with the numbering of its own full document.  The complete native source of the article as distributed in the arXiv e-print for this exact version is bundled unchanged as \texttt{sources/204678/source0.tex}.
\begin{lemma}
\label{lemma1}
Let $X$ be the dynamical vector field of a classical mechanical system on a smooth phase space $M$. Let $f_1,\ldots,f_k\in C^\infty(M)$ satisfy
\begin{equation}
\label{eqPI}
X(f_i)=a_i^j f_j,\qquad i=1,\ldots,k,
\end{equation}
where summation over the repeated index $j$ is understood and $a_i^j\in C^\infty(M)$. Suppose that $f_1,\ldots,f_k$ are functionally independent on
\begin{equation*}
\label{Mf}
M_f=\{x\in M:\ f_1(x)=0,\ldots,f_k(x)=0\}.
\end{equation*}
Then $M_f$ is a dynamically invariant embedded submanifold of $M$ of codimension $k$.
\end{lemma}

\begin{theorem}
\label{maintheorem}
Let $f_1,\ldots,f_k\in C^\infty(M)$, $k<n$, be a system of particular integrals for $(M,\omega,H)$, i.e., 
\begin{equation*}
\{f_i,H\}=a_{i}^j f_j,
\qquad i=1,\ldots,k,
\end{equation*}
for suitable smooth functions $a_i{}^j$; such that they are in particular involution, i.e.,
\begin{equation}
\{f_i,f_j\}=c_{ij}{}^\ell f_\ell,
\qquad i,j=1,\ldots,k,
\label{eq:weak-involution}
\end{equation}
where the $c_{ij}{}^\ell$ are smooth functions. 

Assume that $f_1,\ldots,f_k$ are functionally independent on the level set
\begin{equation*}
M_f=\{x\in M:\ f_1(x)=0,\ldots,f_k(x)=0\}.
\end{equation*}

Then, $M_f$ is a presymplectic manifold of dimension $2n-k$ with the presymplectic form $\Omega=\iota^{*}\omega$, where $\iota\colon M_f \hookrightarrow M$ is the inclusion map, its characteristic distribution is
\begin{equation*}
\ker\Omega=\operatorname{span}\bigl\{X_{f_1}|_{M_f},\ldots,X_{f_k}|_{M_f}\bigr\};
\end{equation*}
and the restricted Hamiltonian $H_{M_f}=H|_{M_f}$ is constant along the characteristic leaves.

If in addition, the characteristic foliation $\ker(\Omega)$ is simple and its leaf space
\[
  \overline{M}=M_f/\ker(\Omega)
\]
is a smooth manifold for which the quotient map
\(\pi:M_f\to \overline{M}\) is a surjective submersion, then there exists a unique symplectic form $\overline{\omega}$ and a unique function
$\overline{H}\in C^\infty(\overline{M})$ such that
\begin{equation}
  \pi^*\overline{\omega}=\Omega,
  \qquad
  \pi^*\overline{H}=H_{M_f}.
  \label{eq:reduced-form-and-Hamiltonian}
\end{equation}
Moreover, $X_H|_{M_f}$ is projectable and
\begin{equation}
  \pi_*(X_H|_{M_f})=X_{\overline{H}},
  \qquad
  \iota_{X_{\overline{H}}}\overline{\omega}
  =\mathrm d\overline{H}.
  \label{eq:projected-Hamiltonian-vector-field}
\end{equation}
In particular, $\overline{M}$ has dimension $2(n-k)$.
\end{theorem}

\begin{proof}
Let
\[
F=(f_1,\ldots,f_k)\colon M\longrightarrow\mathbb{R}^k.
\]
By the functional-independence hypothesis along $M_f$, we have that $0$ is a regular value of $F$, so, from Lemma \ref{lemma1} we have that $M_f$ is a dynamically invariant submanifold of dimension $2n-k$.

Since $0$ is a regular value of $F$, at every $x\in M_f$ one has
\[
  T_xM_f=\bigcap_{i=1}^k\ker(df_i(x))
\]
and therefore
\begin{equation*}
(T_xM_f)^\omega=\operatorname{span}\bigl\{X_{f_1}(x),\ldots,X_{f_k}(x)\bigr\}.
\label{eq:symplectic-orthogonal}
\end{equation*}
For every $i,j$, we have
\[
  X_{f_i}(f_j)|_{M_f}=\{f_j,f_i\}|_{M_f}=0
\]
by \eqref{eq:weak-involution}; so, each $X_{f_i}$ is tangent to $M_f$, consequently, $M_f$ is coisotropic. Furthermore,
\[
  \ker\Omega_x=T_xM_f\cap(T_xM_f)^\omega=(T_xM_f)^\omega.
\]
The independence of the $df_i$ implies that the vectors $X_{f_i}|_{M_f}$ are independent, so $\ker\Omega$ has constant rank $k$.  Since $\dim M_f=2n-k$, it follows that $rank\Omega=2n-2k$.

For every characteristic vector field $X_{f_i}|_{M_f}$, we have
\[
dH_{M_f}(X_{f_i}|_{M_f})
=X_{f_i}(H)|_{M_f}
=\{H,f_i\}|_{M_f}=0.
\]
Therefore, $H_{M_f}$ is constant on the leaves.

If the quotient is smooth, then there is a unique function $\overline{H}$ satisfying the second identity in
\eqref{eq:reduced-form-and-Hamiltonian}.

The form $\Omega$ is horizontal with respect to $\ker(\Omega)$ by definition of the characteristic distribution.  It is also invariant along characteristic vector fields: for $Z\in\ker\Omega$,
\[
\mathcal L_Z\Omega = d(\iota_Z\Omega)+\iota_Z d\Omega=0.
\]
Hence, $\Omega$ is basic and descends to the unique form
$\overline{\omega}$ in \eqref{eq:reduced-form-and-Hamiltonian}.  This form is closed and nondegenerate, and is therefore symplectic.

Finally, because $X_H$ is tangent to $M_f$,
\[
\iota_{X_H|_{M_f}}\Omega= dH_{M_f}.
\]
The flow of $X_H|_{M_f}$ preserves $\Omega$ and its kernel, so $X_H|_{M_f}$ is projectable. Pulling back the defining equation for its projection gives \eqref{eq:projected-Hamiltonian-vector-field}.  Uniqueness follows from the nondegeneracy of $\overline{\omega}$.
\end{proof}

\end{document}
